%% ============================================================================
%% Plantilla de artículo académico - Proyectos aplicados
%% Universidad Santo Tomás - Facultad de Estadística
%% Cursos: Teoría de Respuesta al Ítem, Modelación de Ecuaciones Estructurales y
%%         Muestreo Avanzado
%%
%% Compilación:  pdflatex articulo -> bibtex articulo -> pdflatex -> pdflatex
%% Las referencias se escriben en el archivo referencias.bib. Si ese archivo no existe,
%% la plantilla lo crea con algunos ejemplos (entorno filecontents del preámbulo); si ya
%% existe, no lo sobrescribe. Edite referencias.bib directamente.
%%
%% Los textos en gris son instrucciones: reemplácelos por el contenido del artículo y
%% elimine el comando \guia cuando termine.
%% ============================================================================
\documentclass[11pt,letterpaper]{article}

%% ---------- Paquetes --------------------------------------------------------
\usepackage[utf8]{inputenc}
\usepackage[T1]{fontenc}
\usepackage[spanish,es-noshorthands,es-tabla]{babel}
\usepackage{lmodern}
\usepackage{microtype}
\usepackage[margin=2.5cm]{geometry}
\usepackage{amsmath,amssymb,bm}
\usepackage{graphicx}
\usepackage{booktabs,tabularx,array}
\usepackage[font=small,labelfont=bf]{caption}
\usepackage{subcaption}
\usepackage[dvipsnames]{xcolor}
\usepackage{enumitem}
\usepackage{authblk}
\usepackage{setspace}
\usepackage{tikz}
\usepackage[authoryear,round]{natbib}
\usepackage{xurl}
\usepackage[colorlinks=true,linkcolor=MidnightBlue,citecolor=MidnightBlue,
            urlcolor=MidnightBlue]{hyperref}

\onehalfspacing
\setlength{\emergencystretch}{3em}
\graphicspath{{figuras/}}

%% Texto de instrucciones (bórrelo al terminar)
\newcommand{\guia}[1]{{\color{gray}\itshape #1}}

%% Palabras clave
\newcommand{\palabrasclave}[1]{\par\medskip\noindent\textbf{Palabras clave:} #1}
\newcommand{\keywords}[1]{\par\medskip\noindent\textbf{Keywords:} #1}

%% ---------- Bibliografía de ejemplo ----------------------------------------
\begin{filecontents*}{referencias.bib}
@book{rcore2025,
  author    = {{R Core Team}},
  title     = {R: A Language and Environment for Statistical Computing},
  publisher = {R Foundation for Statistical Computing},
  address   = {Viena},
  year      = {2025},
  url       = {https://www.R-project.org/}
}
@misc{dane2018cnpv,
  author       = {{Departamento Administrativo Nacional de Estadística (DANE)}},
  title        = {Censo Nacional de Población y Vivienda 2018: microdatos anonimizados},
  year         = {2019},
  howpublished = {Archivo Nacional de Datos. \url{https://microdatos.dane.gov.co/index.php/catalog/643}},
  note         = {Consultado el 31 de octubre de 2026}
}
@book{deayala2022,
  author    = {de Ayala, R. J.},
  title     = {The Theory and Practice of Item Response Theory},
  edition   = {2},
  publisher = {Guilford Press},
  address   = {Nueva York},
  year      = {2022}
}
@article{rosseel2012,
  author  = {Rosseel, Yves},
  title   = {lavaan: An {R} Package for Structural Equation Modeling},
  journal = {Journal of Statistical Software},
  year    = {2012},
  volume  = {48},
  number  = {2},
  pages   = {1--36}
}
@book{sarndal1992,
  author    = {Särndal, Carl-Erik and Swensson, Bengt and Wretman, Jan},
  title     = {Model Assisted Survey Sampling},
  publisher = {Springer},
  address   = {Nueva York},
  year      = {1992}
}
\end{filecontents*}

%% ---------- Título y autores -------------------------------------------------
\title{\textbf{Título informativo del artículo: constructo o población, datos y método}}

\author[1]{Nombre Apellido}
\author[1]{Nombre Apellido}
\author[1]{Nombre Apellido}
\affil[1]{Facultad de Estadística, Universidad Santo Tomás, Bogotá, Colombia}
\date{Curso: \underline{\hspace{5cm}} \quad Fecha: \today}

%% ============================================================================
\begin{document}

\maketitle

\begin{abstract}
\noindent\guia{Resumen estructurado de 150 a 250 palabras, en un solo párrafo, sin citas ni
fórmulas. Responda en este orden: (1) objetivo; (2) datos: fuente, población y tamaño;
(3) métodos principales; (4) resultados con las cifras más importantes; (5) conclusión e
implicación.}
\palabrasclave{\guia{tres a cinco términos, separados por punto y coma}}
\end{abstract}

\renewcommand{\abstractname}{Abstract}
\begin{abstract}
\noindent\guia{English version of the abstract.}
\keywords{\guia{three to five terms}}
\end{abstract}

%% ----------------------------------------------------------------------------
\section{Introducción}
\label{sec:introduccion}

\guia{Uno a dos páginas. Presente el problema y su relevancia en el contexto colombiano;
resuma lo que se sabe (con citas) y el vacío que aborda el estudio; cierre con el objetivo
o las preguntas de investigación y una frase que describa la organización del artículo.}

Ejemplo de cita en el texto: el análisis se realizó en R \citep{rcore2025}, con datos del
\citet{dane2018cnpv}.

%% ----------------------------------------------------------------------------
\section{Marco teórico}
\label{sec:marco}

\guia{Presente los conceptos y modelos estadísticos que usa el estudio, con las ecuaciones
necesarias y sus referencias. En SEM, incluya aquí los antecedentes de cada relación y las
hipótesis numeradas. No copie definiciones de libro de texto que el lector no necesita.}

Las ecuaciones se numeran y se citan en el texto, como la ecuación~\eqref{eq:ejemplo}:
\begin{equation}
P(X_{ij}=1\mid\theta_i)=\frac{\exp\{a_j(\theta_i-b_j)\}}{1+\exp\{a_j(\theta_i-b_j)\}} .
\label{eq:ejemplo}
\end{equation}

%% ----------------------------------------------------------------------------
\section{Datos}
\label{sec:datos}

\guia{Fuente (entidad, operación, año, URL), población y unidades, variables o ítems usados
(tabla), procesamiento y síntesis de la calidad de los datos obtenida en la Fase~1.}

\begin{table}[ht]
\centering
\caption{Ejemplo de tabla: variables o ítems utilizados.}
\label{tab:variables}
\small
\begin{tabularx}{\textwidth}{@{}lXl@{}}
\toprule
\textbf{Código} & \textbf{Descripción o enunciado} & \textbf{Escala} \\
\midrule
x1 & Descripción de la variable o enunciado del ítem & Ordinal (1--5) \\
x2 & Descripción de la variable o enunciado del ítem & Dicotómica (0/1) \\
\bottomrule
\end{tabularx}

\smallskip
\parbox{\textwidth}{\footnotesize\textit{Fuente:} elaboración propia con datos de
\citet{dane2018cnpv}.}
\end{table}

%% ----------------------------------------------------------------------------
\section{Métodos}
\label{sec:metodos}

\guia{Describa lo que se hizo con detalle suficiente para que otro estadístico pueda
repetirlo: modelos o diseños, estimación, criterios de decisión (con sus referencias),
tratamiento de faltantes y del diseño muestral, y software con versiones (R y paquetes).}

%% ----------------------------------------------------------------------------
\section{Resultados}
\label{sec:resultados}

\guia{Presente los resultados en el orden de los métodos, con tablas y figuras comentadas en
el texto. Reporte estimaciones con su incertidumbre (error estándar o intervalo). Deje la
interpretación amplia para la discusión.}

\begin{figure}[ht]
\centering
%% Reemplace el recuadro por la figura generada en R, por ejemplo:
%% \includegraphics[width=0.75\textwidth]{curva_informacion.pdf}
\begin{tikzpicture}
\draw[gray,dashed] (0,0) rectangle (10,4.5);
\node[gray,align=center] at (5,2.25) {Inserte aquí la figura con\\
\texttt{\textbackslash includegraphics\{archivo.pdf\}}};
\end{tikzpicture}
\caption{Título descriptivo de la figura.}
\label{fig:ejemplo}

\smallskip
\parbox{\textwidth}{\centering\footnotesize\textit{Fuente:} elaboración propia.}
\end{figure}

La Figura~\ref{fig:ejemplo} y la Tabla~\ref{tab:variables} se citan con \texttt{\textbackslash
ref}.

%% ----------------------------------------------------------------------------
\section{Discusión}
\label{sec:discusion}

\guia{Interprete los resultados a la luz de las preguntas y de la literatura; señale las
implicaciones prácticas y las limitaciones (datos, diseño, supuestos del modelo); proponga
trabajo futuro.}

%% ----------------------------------------------------------------------------
\section{Conclusiones}
\label{sec:conclusiones}

\guia{Uno o dos párrafos que respondan directamente el objetivo, sin repetir cifras ya
discutidas ni introducir resultados nuevos.}

%% ----------------------------------------------------------------------------
\section*{Disponibilidad de datos y código}

\guia{Indique la fuente y la URL de los datos y dónde se encuentran los scripts de R que
reproducen todos los resultados (por ejemplo, el archivo comprimido entregado o un
repositorio).}

%% ----------------------------------------------------------------------------
%% Referencias que no se citan en el texto pero deben aparecer (opcional):
\nocite{deayala2022,rosseel2012,sarndal1992}

\bibliographystyle{apalike}
\bibliography{referencias}

%% ----------------------------------------------------------------------------
\appendix
\section{Material complementario}
\label{app:complementario}

\guia{Tablas extensas, resultados complementarios y sintaxis del modelo.}

\end{document}
